% Options for packages loaded elsewhere
\PassOptionsToPackage{unicode}{hyperref}
\PassOptionsToPackage{hyphens}{url}
%
\documentclass[
]{article}
\usepackage{amsmath,amssymb}
\usepackage{iftex}
\ifPDFTeX
  \usepackage[T1]{fontenc}
  \usepackage[utf8]{inputenc}
  \usepackage{textcomp} % provide euro and other symbols
\else % if luatex or xetex
  \usepackage{unicode-math} % this also loads fontspec
  \defaultfontfeatures{Scale=MatchLowercase}
  \defaultfontfeatures[\rmfamily]{Ligatures=TeX,Scale=1}
\fi
\usepackage{lmodern}
\ifPDFTeX\else
  % xetex/luatex font selection
\fi
% Use upquote if available, for straight quotes in verbatim environments
\IfFileExists{upquote.sty}{\usepackage{upquote}}{}
\IfFileExists{microtype.sty}{% use microtype if available
  \usepackage[]{microtype}
  \UseMicrotypeSet[protrusion]{basicmath} % disable protrusion for tt fonts
}{}
\makeatletter
\@ifundefined{KOMAClassName}{% if non-KOMA class
  \IfFileExists{parskip.sty}{%
    \usepackage{parskip}
  }{% else
    \setlength{\parindent}{0pt}
    \setlength{\parskip}{6pt plus 2pt minus 1pt}}
}{% if KOMA class
  \KOMAoptions{parskip=half}}
\makeatother
\usepackage{xcolor}
\usepackage{color}
\usepackage{fancyvrb}
\newcommand{\VerbBar}{|}
\newcommand{\VERB}{\Verb[commandchars=\\\{\}]}
\DefineVerbatimEnvironment{Highlighting}{Verbatim}{commandchars=\\\{\}}
% Add ',fontsize=\small' for more characters per line
\newenvironment{Shaded}{}{}
\newcommand{\AlertTok}[1]{\textcolor[rgb]{1.00,0.00,0.00}{\textbf{#1}}}
\newcommand{\AnnotationTok}[1]{\textcolor[rgb]{0.38,0.63,0.69}{\textbf{\textit{#1}}}}
\newcommand{\AttributeTok}[1]{\textcolor[rgb]{0.49,0.56,0.16}{#1}}
\newcommand{\BaseNTok}[1]{\textcolor[rgb]{0.25,0.63,0.44}{#1}}
\newcommand{\BuiltInTok}[1]{\textcolor[rgb]{0.00,0.50,0.00}{#1}}
\newcommand{\CharTok}[1]{\textcolor[rgb]{0.25,0.44,0.63}{#1}}
\newcommand{\CommentTok}[1]{\textcolor[rgb]{0.38,0.63,0.69}{\textit{#1}}}
\newcommand{\CommentVarTok}[1]{\textcolor[rgb]{0.38,0.63,0.69}{\textbf{\textit{#1}}}}
\newcommand{\ConstantTok}[1]{\textcolor[rgb]{0.53,0.00,0.00}{#1}}
\newcommand{\ControlFlowTok}[1]{\textcolor[rgb]{0.00,0.44,0.13}{\textbf{#1}}}
\newcommand{\DataTypeTok}[1]{\textcolor[rgb]{0.56,0.13,0.00}{#1}}
\newcommand{\DecValTok}[1]{\textcolor[rgb]{0.25,0.63,0.44}{#1}}
\newcommand{\DocumentationTok}[1]{\textcolor[rgb]{0.73,0.13,0.13}{\textit{#1}}}
\newcommand{\ErrorTok}[1]{\textcolor[rgb]{1.00,0.00,0.00}{\textbf{#1}}}
\newcommand{\ExtensionTok}[1]{#1}
\newcommand{\FloatTok}[1]{\textcolor[rgb]{0.25,0.63,0.44}{#1}}
\newcommand{\FunctionTok}[1]{\textcolor[rgb]{0.02,0.16,0.49}{#1}}
\newcommand{\ImportTok}[1]{\textcolor[rgb]{0.00,0.50,0.00}{\textbf{#1}}}
\newcommand{\InformationTok}[1]{\textcolor[rgb]{0.38,0.63,0.69}{\textbf{\textit{#1}}}}
\newcommand{\KeywordTok}[1]{\textcolor[rgb]{0.00,0.44,0.13}{\textbf{#1}}}
\newcommand{\NormalTok}[1]{#1}
\newcommand{\OperatorTok}[1]{\textcolor[rgb]{0.40,0.40,0.40}{#1}}
\newcommand{\OtherTok}[1]{\textcolor[rgb]{0.00,0.44,0.13}{#1}}
\newcommand{\PreprocessorTok}[1]{\textcolor[rgb]{0.74,0.48,0.00}{#1}}
\newcommand{\RegionMarkerTok}[1]{#1}
\newcommand{\SpecialCharTok}[1]{\textcolor[rgb]{0.25,0.44,0.63}{#1}}
\newcommand{\SpecialStringTok}[1]{\textcolor[rgb]{0.73,0.40,0.53}{#1}}
\newcommand{\StringTok}[1]{\textcolor[rgb]{0.25,0.44,0.63}{#1}}
\newcommand{\VariableTok}[1]{\textcolor[rgb]{0.10,0.09,0.49}{#1}}
\newcommand{\VerbatimStringTok}[1]{\textcolor[rgb]{0.25,0.44,0.63}{#1}}
\newcommand{\WarningTok}[1]{\textcolor[rgb]{0.38,0.63,0.69}{\textbf{\textit{#1}}}}
\usepackage{longtable,booktabs,array}
\usepackage{calc} % for calculating minipage widths
% Correct order of tables after \paragraph or \subparagraph
\usepackage{etoolbox}
\makeatletter
\patchcmd\longtable{\par}{\if@noskipsec\mbox{}\fi\par}{}{}
\makeatother
% Allow footnotes in longtable head/foot
\IfFileExists{footnotehyper.sty}{\usepackage{footnotehyper}}{\usepackage{footnote}}
\makesavenoteenv{longtable}
\setlength{\emergencystretch}{3em} % prevent overfull lines
\providecommand{\tightlist}{%
  \setlength{\itemsep}{0pt}\setlength{\parskip}{0pt}}
\setcounter{secnumdepth}{-\maxdimen} % remove section numbering
\ifLuaTeX
  \usepackage{selnolig}  % disable illegal ligatures
\fi
\IfFileExists{bookmark.sty}{\usepackage{bookmark}}{\usepackage{hyperref}}
\IfFileExists{xurl.sty}{\usepackage{xurl}}{} % add URL line breaks if available
\urlstyle{same}
\hypersetup{
  hidelinks,
  pdfcreator={LaTeX via pandoc}}

\author{}
\date{}

\begin{document}

\hypertarget{suffering-aware-neural-networks-on-amd-alveo-u250-deploying-the-exit-audit-flop-metering-kernel}{%
\section{Suffering-Aware Neural Networks on AMD Alveo U250: Deploying
the Exit-Audit / FLOP-Metering
Kernel}\label{suffering-aware-neural-networks-on-amd-alveo-u250-deploying-the-exit-audit-flop-metering-kernel}}

\textbf{Status:} \texttt{DRAFT} --- every empirical claim marked
\texttt{MEASURED} below is backed by a measured artifact in the Sounio
repository; claims marked \texttt{ESTIMATE} or \texttt{(in\ progress)}
are explicitly noted as such. Human measurement audit is the authority;
LLM reviews are recorded in Section7 only as an AI-assist disclosure.
\textbf{Date:} 2026-08-04 \textbf{Orthography:} EN-US \textbf{Companion
spec:}
\texttt{docs/research/san\_imagenet\_fpga\_dl380\_spec\_2026-08-02.md}
\textbf{Companion contract:}
\texttt{scripts/research/san\_imagenet\_fpga\_dl380.py} \textbf{Target
venue:} arXiv \texttt{cs.LG} / \texttt{cs.AR}; then IEEE FPL or a
systems note

\begin{center}\rule{0.5\linewidth}{0.5pt}\end{center}

\hypertarget{abstract}{%
\subsection{Abstract}\label{abstract}}

We report the first measured deployment of the SAN exit-audit /
FLOP-metering kernel on an AMD Alveo U250 FPGA, together with a GPU
training study of SAN-ResNet-50, SAN-ViT-small/d384, and a tiny decoder
LM (SAN-GPT-small). Suffering-aware neural networks (SANs) are
early-exit architectures whose training freezes as soon as a held-out
feasibility target is met, eliminating gratuitous computation. We
offload only the inference-time catastrophe-scan and FLOP-metering path
to the FPGA; the trunk runs on the host CPU/GPU. On the U250 the scan
kernel runs at \textbf{511 Msamples/s sustained} on a 1.2M-sample stress
cohort (\textbf{94.5\%} of the 540.8 Msamples/s bus-limited theoretical
peak at the achieved 135.2 MHz clock) and consumes approximately
\textbf{3.3 nJ/sample} incremental board-level energy (rounded to the
precision the on-card power sensor supports). The kernel is bit-exact
against the control-VM golden model on synthetic cohorts, the stress
cohort, and real photographs from ImageNette2-160.

On NVIDIA GPUs (A5000 / RTX 4000 Ada) SAN training reduces the
\emph{metered-MAC} integrated machine burden by \textbf{40.7\%}
(ResNet-50), \textbf{50.4\%} (ViT-small/d384), and \textbf{52.2\%}
(SAN-GPT-small) relative to dense training in a small-subset,
fast-convergence regime (validation accuracies 0.39, 0.26, 0.17, chosen
to demonstrate the freeze-on-green mechanism, not to claim competitive
task accuracy). ResNet-50 also shows a measured 1.08x wall-time latency
speedup on CIFAR-10, though this margin is small and task-dependent. We
disclose the limits honestly: the energy number is board-level, not
rack-level; full ImageNet-1k is unavailable, so ImageNette2-160 is the
real-image proxy; and ResNet-50 shows a disclosed patient-channel
tradeoff while the other families satisfy the stricter L5 clause in this
run. All artifacts, measurements, and reproduction scripts are in the
repository.

\begin{center}\rule{0.5\linewidth}{0.5pt}\end{center}

\hypertarget{introduction}{%
\subsection{1 Introduction}\label{introduction}}

\hypertarget{the-problem-computation-we-do-not-need}{%
\subsubsection{1.1 The problem: computation we do not
need}\label{the-problem-computation-we-do-not-need}}

Deep learning at scale pays for every layer it runs, including layers
that do not change the answer. Early-exit networks {[}1,2,3{]} reduce
this cost by allowing samples to leave the network at intermediate
classifiers when they are ``confident enough.'' Most early-exit work,
however, optimizes for speedup or FLOP reduction alone, and trains for a
fixed budget; it therefore continues to spend computation after the
model has already become good enough on the task it was asked to solve.

A suffering-aware neural network (SAN) {[}4{]} inverts the framing. It
defines two channels of suffering:

\begin{itemize}
\tightlist
\item
  \textbf{Patient suffering}: the cost of prediction errors under an
  asymmetric harm structure (e.g., missing a hazard class is worse than
  a false alarm).
\item
  \textbf{Machine suffering}: the computation actually executed, metered
  exactly.
\end{itemize}

Training proceeds only until a held-out feasibility target is reached
(e.g., validation accuracy \textgreater= tau). Once the target is met,
training freezes \emph{immediately}. The result is a strict separation
of machine suffering into a \emph{necessary} part (up to the freeze
epoch) and a \emph{gratuitous} part (everything after), with gratuitous
suffering forced to zero. At inference, early exits further reduce
per-sample machine suffering whenever a sample can exit early.

This paper asks: can this idea be deployed as a real system? We answer
with a complete pipeline --- training on GPU, confidence extraction, and
a bit-exact FPGA catastrophe-scan / FLOP-metering kernel on an AMD Alveo
U250 --- measured end to end.

\hypertarget{contributions}{%
\subsubsection{1.2 Contributions}\label{contributions}}

\begin{longtable}[]{@{}
  >{\raggedright\arraybackslash}p{(\columnwidth - 6\tabcolsep) * \real{0.2500}}
  >{\raggedright\arraybackslash}p{(\columnwidth - 6\tabcolsep) * \real{0.2500}}
  >{\raggedright\arraybackslash}p{(\columnwidth - 6\tabcolsep) * \real{0.2500}}
  >{\raggedright\arraybackslash}p{(\columnwidth - 6\tabcolsep) * \real{0.2500}}@{}}
\toprule\noalign{}
\begin{minipage}[b]{\linewidth}\raggedright
\#
\end{minipage} & \begin{minipage}[b]{\linewidth}\raggedright
Contribution
\end{minipage} & \begin{minipage}[b]{\linewidth}\raggedright
Evidence
\end{minipage} & \begin{minipage}[b]{\linewidth}\raggedright
Status
\end{minipage} \\
\midrule\noalign{}
\endhead
\bottomrule\noalign{}
\endlastfoot
1 & A SAN training harness for ResNet-50, ViT-small/d384, and a tiny
decoder LM on real data, with metered-MAC accounting and freeze-on-green
& \texttt{scripts/research/suffering\_aware\_large\_architecture.py} &
\texttt{MEASURED} on GPU (small subset) \\
2 & An FPGA exit-audit / FLOP-meter kernel with integer semantics, no
multipliers/DSPs, and one sample/cycle/PE throughput &
\texttt{hardware/fpga/u250\_catastrophe\_scan/krnl\_san\_scan.cpp} &
\texttt{MEASURED} on U250 \\
3 & On-target U250 benchmark: 511 Msamples/s on 1.2M stress cohort,
\textasciitilde3.3 nJ/sample board-level energy, bit-exact against
golden model & \texttt{host\_san\_scan}, \texttt{host\_san\_scan\_bench}
& \texttt{MEASURED} \\
4 & GPU training study: 40.7--52.2\% metered-MAC savings in a
fast-convergence, small-subset regime; CIFAR-100 is infeasible under the
same budget & Slurm jobs 8584/8588/8591, 8599--8607, 8611--8612 &
\texttt{MEASURED\ /\ MIXED} \\
5 & Real-image kernel validation on ImageNette2-160, bit-exact on the
U250 & \texttt{train\_san\_imagenette.py} + \texttt{host\_san\_scan} &
\texttt{MEASURED} \\
6 & Honest accounting of limits: small-subset accuracy, partial meter
convention, ImageNet-1k unavailable, board-level power, ResNet-50
patient-channel tradeoff & Section5, Section6 & \texttt{DECLARED} \\
\end{longtable}

The primary systems contribution is the measured bridge between the SAN
learning rule and FPGA deployment: the kernel is not an approximate
accelerator of a float model; it \emph{is} the integer specification of
the exit-audit and FLOP-metering path, and the deployment is sound
exactly when the card reproduces that specification. The trunk stays on
the host; the card decides, counts, and meters.

\hypertarget{prior-work}{%
\subsubsection{1.3 Prior work}\label{prior-work}}

Early-exit networks were introduced by BranchyNet {[}1{]} and
Shallow-Deep Networks {[}2{]}, with later work on confidence thresholds
{[}3{]}, dynamic inference, and hardware-aware early exit. Most of this
work treats the threshold as a hyperparameter tuned for
accuracy-efficiency tradeoffs and trains for a fixed budget.

Constrained learning {[}5,6,7{]} provides generalization guarantees for
ERM with explicit constraints, but does not address inference-time
compute reduction or the freeze-on-green training rule. Mercyful /
suffering-aware learning {[}4{]} introduces the two-channel suffering
framework and the anti-Goodhart selection rule; the present paper is the
first systems study that maps that framework to a measured FPGA
deployment.

On the FPGA side, early-exit accelerators have been proposed for CNNs
and transformers, typically using custom datapaths that map part of the
model into hardware. Our kernel takes the opposite approach: the trunk
stays on the host (or GPU), and only the cheap decision/metering path
runs on the FPGA. This keeps the bitstream small, architecture-agnostic
(the stage-cost LUT is loaded by the host), and verifiable by direct
comparison to a golden model.

\begin{center}\rule{0.5\linewidth}{0.5pt}\end{center}

\hypertarget{background-suffering-aware-neural-networks}{%
\subsection{2 Background: Suffering-Aware Neural
Networks}\label{background-suffering-aware-neural-networks}}

\hypertarget{the-two-channels}{%
\subsubsection{2.1 The two channels}\label{the-two-channels}}

A SAN is defined by:

\begin{itemize}
\tightlist
\item
  A trunk with intermediate exit heads.
\item
  A per-sample exit rule: the first head whose confidence \textgreater=
  Delta decides the output; if no head is confident, the final head
  decides.
\item
  A feasibility target tau on held-out accuracy.
\item
  A training rule: train until tau is reached, then freeze.
\item
  A compassion grid over (patient weight, machine weight) that selects
  among feasible checkpoints.
\end{itemize}

\textbf{Patient suffering} is the mean harm of the model's predictions
under a cost matrix C, where C{[}y, y\_hat{]} encodes the cost of
predicting y\_hat when the true label is y. For CIFAR-10 we use the
parent line's hazard structure: class 9 (truck) is the hazard class;
missing the hazard costs 5, a false hazard alarm costs 2, any other
confusion costs 1. For SAN-GPT-small the hazard tokens are corpus
negation tokens.

\textbf{Machine suffering} is the \emph{metered-MAC} count of executed
operations. We follow the convention used throughout the SAN line: MACs
x 2 = FLOPs; training backward pass = 2x forward; biases, norms,
activations, softmax, residual adds, and pooling are unmetered. The
convention is identical for every architecture and baseline, so all
reported savings are relative under the same partial accounting. It is a
proxy for computational burden, not a claim that every host- side
operation has been charged.

\hypertarget{anti-goodhart-gating}{%
\subsubsection{2.2 Anti-Goodhart gating}\label{anti-goodhart-gating}}

A standard risk when optimizing a proxy (FLOPs) is that the model learns
to satisfy the proxy while failing the real task {[}8{]}. SAN prevents
this by selection: the compassion grid may only choose checkpoints that
are feasible (accuracy \textgreater= tau). If no checkpoint is feasible,
the gate returns \texttt{NO\_FEASIBLE} instead of silently degrading
accuracy. This is an executable property checked in the companion
contract.

\hypertarget{freeze-on-green}{%
\subsubsection{2.3 Freeze-on-green}\label{freeze-on-green}}

Let t* be the first epoch at which the held-out accuracy reaches tau.
SAN stops at t\emph{. The dense baseline trains for the full budget B;
the EarlyStop baseline also stops at t} but lacks early-exit layers.
SAN's integrated machine suffering decomposes as:

S\_SAN = S\_necessary(t*) + S\_gratuitous

with S\_gratuitous = 0 by construction. The dense baseline has
S\_gratuitous \textgreater{} 0; EarlyStop has S\_gratuitous = 0 but
lacks inference-time savings.

\begin{center}\rule{0.5\linewidth}{0.5pt}\end{center}

\hypertarget{methods}{%
\subsection{3 Methods}\label{methods}}

\hypertarget{architectures}{%
\subsubsection{3.1 Architectures}\label{architectures}}

\textbf{SAN-ResNet-50.} A CIFAR-variant ResNet-50 with bottleneck blocks
(3,4,6,3), stage widths 256/512/1024/2048, and one early-exit head after
each stage plus the final head. Real-scale MACs per stage are taken from
the published architecture (Section3.1 of the companion spec).

\textbf{SAN-ViT-small/d384.} Patch 4x4 -\textgreater{} 64 tokens + CLS,
d=384, 12 blocks, 6 heads, MLP ratio 4. One CLS exit head per block.
This is a small attention proxy for the CIFAR pilot, not ViT-Large. The
FPGA kernel's stage-cost LUT can be reloaded with the published ViT-L/16
constants (61.55 GMAC) for a real-scale deployment.

\textbf{SAN-GPT-small.} A tiny decoder-only transformer (d=384, 10
blocks, 6 heads, causal mask, vocab 2000) trained on next-token
prediction over the repository's own research documentation corpus. Exit
heads score the last G=4 positions; the gate confidence is the mean
max-probability over those positions. This is an internal-corpus pilot,
not a general language-modeling result.

All three families are compared against two baselines:

\begin{itemize}
\tightlist
\item
  \textbf{Dense}: the same trunk trained for the full budget.
\item
  \textbf{EarlyStop}: the same trunk with SAN's stop rule but no exit
  heads.
\end{itemize}

\hypertarget{fpga-kernel}{%
\subsubsection{3.2 FPGA kernel}\label{fpga-kernel}}

The kernel \texttt{krnl\_san\_scan.cpp} is a pure integer exit-audit /
FLOP-meter. It does \emph{not} run the trunk; the trunk produces
per-sample confidence vectors on the host CPU/GPU and the kernel
receives them as Q0.15 integers. Per lane it does:

\begin{enumerate}
\def\labelenumi{\arabic{enumi}.}
\tightlist
\item
  Read a packed 512-bit beat containing four 128-bit samples; each
  sample holds up to seven 15-bit Q0.15 confidence fields.
\item
  Compare each field to the integer threshold q\_Delta = floor(Delta .
  2\^{}15).
\item
  Return the index of the first field \textgreater= q\_Delta (priority
  encoder).
\item
  If none, return the final index (catastrophe).
\item
  Look up the stage-cost LUT at the exit index and accumulate into a
  64-bit counter.
\end{enumerate}

The kernel contains \textbf{no floating-point, no softmax, no
multipliers, no DSPs}. It is bus-limited: 512 bits/cycle x 4
samples/beat x 135.2 MHz = 540.8 Msamples/s theoretical peak. The LUT is
loaded by the host, so the same bitstream serves ResNet-50,
ViT-small/d384, and SAN-GPT-small by reloading the LUT and threshold.

\textbf{Correctness criterion (T3).} The kernel is correct iff it
reproduces the golden integer function. We verify this bit-exactly on
three cohorts: the ResNet validation cohort (5,000 samples, 5 exit
points), the ViT validation cohort (5,000 samples, 7 exit points), and a
1.2M-sample stress cohort (5 exit points, ImageNet-sized). The host
program \texttt{host\_san\_scan} compares card outputs to the control-VM
golden and reports \texttt{HOST\_SAN\_SCAN\_PASS}.

\hypertarget{gpu-training-harness}{%
\subsubsection{3.3 GPU training harness}\label{gpu-training-harness}}

Training runs on the Slurm \texttt{gpu-orangefs} partition (NVIDIA RTX
A5000 / RTX 4000 Ada). The harness is
\texttt{scripts/research/suffering\_aware\_large\_architecture.py},
adapted for CUDA with \texttt{torch.cuda.synchronize()} latency
measurement. Each family shares one trunk initialization, one data
order, and one seed across SAN, Dense, and EarlyStop.

\hypertarget{datasets}{%
\subsubsection{3.4 Datasets}\label{datasets}}

\begin{itemize}
\tightlist
\item
  \textbf{CIFAR-10}: 50,000 train / 10,000 test, 10 classes. Used for
  the main GPU training study.
\item
  \textbf{ImageNette2-160}: 10-class subset of ImageNet, 160 px, real
  photographs. Used for real-image U250 validation. Full ImageNet-1k is
  not available in this environment.
\item
  \textbf{CIFAR-100 / ImageNette2-320 / ImageNet-1k}: ablation /
  larger-proxy datasets; availability and results reported in
  Section4.3--4.4.
\end{itemize}

\begin{center}\rule{0.5\linewidth}{0.5pt}\end{center}

\hypertarget{experimental-results}{%
\subsection{4 Experimental Results}\label{experimental-results}}

\hypertarget{u250-throughput-and-energy}{%
\subsubsection{4.1 U250 throughput and
energy}\label{u250-throughput-and-energy}}

The bitstream was built with Vitis 2025.1 on a separate builder VM and
run on the DL380 U250. Table 1 reports the on-target campaign.

\textbf{Table 1: U250 on-target benchmark campaign.}

\begin{longtable}[]{@{}
  >{\raggedright\arraybackslash}p{(\columnwidth - 10\tabcolsep) * \real{0.1667}}
  >{\raggedright\arraybackslash}p{(\columnwidth - 10\tabcolsep) * \real{0.1667}}
  >{\raggedright\arraybackslash}p{(\columnwidth - 10\tabcolsep) * \real{0.1667}}
  >{\raggedright\arraybackslash}p{(\columnwidth - 10\tabcolsep) * \real{0.1667}}
  >{\raggedright\arraybackslash}p{(\columnwidth - 10\tabcolsep) * \real{0.1667}}
  >{\raggedright\arraybackslash}p{(\columnwidth - 10\tabcolsep) * \real{0.1667}}@{}}
\toprule\noalign{}
\begin{minipage}[b]{\linewidth}\raggedright
dataset
\end{minipage} & \begin{minipage}[b]{\linewidth}\raggedright
n
\end{minipage} & \begin{minipage}[b]{\linewidth}\raggedright
points
\end{minipage} & \begin{minipage}[b]{\linewidth}\raggedright
single-shot (Msamples/s)
\end{minipage} & \begin{minipage}[b]{\linewidth}\raggedright
sustained (Msamples/s)
\end{minipage} & \begin{minipage}[b]{\linewidth}\raggedright
result
\end{minipage} \\
\midrule\noalign{}
\endhead
\bottomrule\noalign{}
\endlastfoot
val\_resnet & 5,000 & 5 & 24.2 & 146.7 &
\texttt{HOST\_SAN\_SCAN\_PASS} \\
val\_vit & 5,000 & 7 & 43.0 & 146.9 & \texttt{HOST\_SAN\_SCAN\_PASS} \\
stress\_1p2M & 1,200,000 & 5 & 481.9 & \textbf{511.0} &
\texttt{HOST\_SAN\_SCAN\_PASS} \\
val\_imagenette & 3,925 & 5 & 24.1 & 122.2 &
\texttt{HOST\_SAN\_SCAN\_PASS} \\
\end{longtable}

The 1.2M stress cohort reaches \textbf{511 Msamples/s sustained},
\textbf{94.5\%} of the 540.8 Msamples/s theoretical peak at the achieved
135.2 MHz clock. This is a best-case stress microbenchmark; smaller
cohorts are enqueue/sync dominated and report lower sustained rates
(Table 1). We report both numbers honestly rather than lead only with
the large-cohort peak.

\textbf{Energy.} Board-level power was measured with
\texttt{xrt-smi\ examine\ -r\ electrical} (1 Hz, 30 s). Idle card:
24.435 W. Under continuous \texttt{host\_san\_scan\_bench} on the 1.2M
cohort: 26.153 W. Incremental draw DeltaP \textasciitilde{} 1.7 W. The
bench processed 15.5436 Gsamples in 30.002 s (aggregate 518.1
Msamples/s), giving approximately \textbf{3.3 nJ/sample} incremental
board-level energy. A repeat with the tiny ImageNette cohort gave load
power below idle because the kernel idles between micro-enqueues; the
stress-cohort number is therefore the honest per-sample energy figure.
We round to two significant figures because the on-card power sensor is
sampled at 1 Hz.

\hypertarget{gpu-training-savings-and-latency}{%
\subsubsection{4.2 GPU training savings and
latency}\label{gpu-training-savings-and-latency}}

Table 2 reports the GPU scale pilot. All numbers are measured on real
CIFAR-10 training with \texttt{SAN\_LARGE\_DEVICE=cuda}, using
stratified subsets of 4,000 train / 1,000 val images and a small epoch
budget (8 for ResNet-50, 10 for ViT-small/d384/SAN-GPT-small). The
feasibility targets tau (0.34, 0.251, 0.165) are intentionally low so
the freeze-on-green rule can be demonstrated quickly; the achieved
accuracies are not competitive with full-dataset CIFAR-10 baselines. The
point of the table is to show that the SAN machinery (meter
conservation, freeze-on-green, early exits, compassion grid) operates
correctly across families, not to claim a new accuracy result.

\textbf{Table 2: GPU training study (CIFAR-10, small subset).}

\begin{longtable}[]{@{}
  >{\raggedright\arraybackslash}p{(\columnwidth - 18\tabcolsep) * \real{0.1000}}
  >{\raggedright\arraybackslash}p{(\columnwidth - 18\tabcolsep) * \real{0.1000}}
  >{\raggedright\arraybackslash}p{(\columnwidth - 18\tabcolsep) * \real{0.1000}}
  >{\raggedright\arraybackslash}p{(\columnwidth - 18\tabcolsep) * \real{0.1000}}
  >{\raggedright\arraybackslash}p{(\columnwidth - 18\tabcolsep) * \real{0.1000}}
  >{\raggedright\arraybackslash}p{(\columnwidth - 18\tabcolsep) * \real{0.1000}}
  >{\raggedright\arraybackslash}p{(\columnwidth - 18\tabcolsep) * \real{0.1000}}
  >{\raggedright\arraybackslash}p{(\columnwidth - 18\tabcolsep) * \real{0.1000}}
  >{\raggedright\arraybackslash}p{(\columnwidth - 18\tabcolsep) * \real{0.1000}}
  >{\raggedright\arraybackslash}p{(\columnwidth - 18\tabcolsep) * \real{0.1000}}@{}}
\toprule\noalign{}
\begin{minipage}[b]{\linewidth}\raggedright
family
\end{minipage} & \begin{minipage}[b]{\linewidth}\raggedright
t*
\end{minipage} & \begin{minipage}[b]{\linewidth}\raggedright
SAN acc@t*
\end{minipage} & \begin{minipage}[b]{\linewidth}\raggedright
S\_m(SAN)
\end{minipage} & \begin{minipage}[b]{\linewidth}\raggedright
S\_m(Dense)
\end{minipage} & \begin{minipage}[b]{\linewidth}\raggedright
saving
\end{minipage} & \begin{minipage}[b]{\linewidth}\raggedright
latency SAN
\end{minipage} & \begin{minipage}[b]{\linewidth}\raggedright
latency Dense
\end{minipage} & \begin{minipage}[b]{\linewidth}\raggedright
speedup
\end{minipage} & \begin{minipage}[b]{\linewidth}\raggedright
verdict
\end{minipage} \\
\midrule\noalign{}
\endhead
\bottomrule\noalign{}
\endlastfoot
ResNet-50 & 4 & 0.390 & 160.1 TMAC & 269.9 TMAC & \textbf{40.7\%} &
0.196 ms & 0.213 ms & 1.08x & L1--L4, L6--L8 PASS; L5 tradeoff \\
ViT-small/d384 & 4 & 0.262 & 183.3 TMAC & 369.3 TMAC & \textbf{50.4\%} &
0.310 ms & 0.308 ms & 0.99x & L1--L4, L7--L8 PASS; L5 PASS, L6 exit-frac
0.03 \\
SAN-GPT-small & 4 & 0.167 & 115.4 TMAC & 241.5 TMAC & \textbf{52.2\%} &
0.343 ms & 0.341 ms & 0.99x & L1--L8 PASS \\
\end{longtable}

\emph{S\_m is the metered-MAC burden (MACx2, backward = 2x forward,
biases/norms/etc. unmetered). Verdicts refer to the companion contract
clauses L1--L8; see Appendix B for clause definitions.}

\textbf{Machine channel.} SAN saves metered-MAC burden in every family.
ResNet-50, with its early-exit-friendly residual stages, also translates
the MAC savings into a small real wall-time speedup on CIFAR-10 (0.196
ms vs 0.213 ms), though the margin is within the noise regime of CUDA
synchronize microbenchmarks. ViT-small/d384 and SAN-GPT-small are
essentially tied with Dense in wall time because the forward is short
enough that dispatch overhead dominates; at larger scale the MAC savings
would likely translate to latency savings as well.

\textbf{Patient channel.} ResNet-50 shows the disclosed tradeoff the
parent line reports at ImageNet scale: SAN patient harm is slightly
higher than EarlyStop's because early stopping alone can freeze on a
luckier epoch. ViT-small/d384 and SAN-GPT-small satisfy the stricter L5
clause (SAN \textless= both baselines) in this run. These are results,
not tuning failures. The cost matrix C is synthetic; no clinical claim
is made.

\hypertarget{threshold-ablation}{%
\subsubsection{4.3 Threshold ablation}\label{threshold-ablation}}

Table 3 reports the Delta sensitivity study on the same CIFAR-10 small
subset (Slurm jobs 8599--8607). The threshold changes which samples exit
early and therefore the metered-MAC burden, but it cannot overcome the
small training budget: even the best configurations remain
\texttt{L\_RED} because the patient-channel clauses (L5/L7 in
particular) are not fully satisfied on this subset. The ablation is
therefore reported as a \emph{sensitivity knob}, not as a tuning recipe
that turns the study green.

\textbf{Table 3: Threshold ablation (CIFAR-10, 4,000 train / 1,000
val).}

\begin{longtable}[]{@{}
  >{\raggedright\arraybackslash}p{(\columnwidth - 16\tabcolsep) * \real{0.1111}}
  >{\raggedright\arraybackslash}p{(\columnwidth - 16\tabcolsep) * \real{0.1111}}
  >{\raggedright\arraybackslash}p{(\columnwidth - 16\tabcolsep) * \real{0.1111}}
  >{\raggedright\arraybackslash}p{(\columnwidth - 16\tabcolsep) * \real{0.1111}}
  >{\raggedright\arraybackslash}p{(\columnwidth - 16\tabcolsep) * \real{0.1111}}
  >{\raggedright\arraybackslash}p{(\columnwidth - 16\tabcolsep) * \real{0.1111}}
  >{\raggedright\arraybackslash}p{(\columnwidth - 16\tabcolsep) * \real{0.1111}}
  >{\raggedright\arraybackslash}p{(\columnwidth - 16\tabcolsep) * \real{0.1111}}
  >{\raggedright\arraybackslash}p{(\columnwidth - 16\tabcolsep) * \real{0.1111}}@{}}
\toprule\noalign{}
\begin{minipage}[b]{\linewidth}\raggedright
family
\end{minipage} & \begin{minipage}[b]{\linewidth}\raggedright
Delta
\end{minipage} & \begin{minipage}[b]{\linewidth}\raggedright
t*
\end{minipage} & \begin{minipage}[b]{\linewidth}\raggedright
SAN acc@t*
\end{minipage} & \begin{minipage}[b]{\linewidth}\raggedright
S\_m(SAN)
\end{minipage} & \begin{minipage}[b]{\linewidth}\raggedright
S\_m(Dense)
\end{minipage} & \begin{minipage}[b]{\linewidth}\raggedright
saving
\end{minipage} & \begin{minipage}[b]{\linewidth}\raggedright
speedup
\end{minipage} & \begin{minipage}[b]{\linewidth}\raggedright
clauses PASS
\end{minipage} \\
\midrule\noalign{}
\endhead
\bottomrule\noalign{}
\endlastfoot
ResNet-50 & 0.35 & --- & 0.338 & 201.7 TMAC & 269.9 TMAC & 25.3\% &
1.54x & 4/8 \\
ResNet-50 & 0.45 & 4 & 0.355 & 149.7 TMAC & 269.9 TMAC & 44.5\% & 1.10x
& 7/8 \\
ResNet-50 & 0.55 & 3 & 0.347 & 130.9 TMAC & 269.9 TMAC & 51.5\% & 1.00x
& 6/8 \\
ResNet-50 & 0.65 & 2 & 0.354 & 100.3 TMAC & 269.9 TMAC & 62.8\% & 0.98x
& 7/8 \\
ResNet-50 & 0.75 & 4 & 0.377 & 167.1 TMAC & 269.9 TMAC & 38.1\% & 0.99x
& 7/8 \\
ViT-small/d384 & 0.25 & --- & 0.187 & 234.6 TMAC & 369.3 TMAC & 36.5\% &
1.83x & 4/8 \\
ViT-small/d384 & 0.35 & 9 & 0.267 & 312.2 TMAC & 369.3 TMAC & 15.5\% &
1.19x & 6/8 \\
ViT-small/d384 & 0.45 & 4 & 0.258 & 183.6 TMAC & 369.3 TMAC & 50.3\% &
0.98x & 6/8 \\
ViT-small/d384 & 0.55 & 4 & 0.257 & 184.2 TMAC & 369.3 TMAC & 50.1\% &
0.99x & 7/8 \\
\end{longtable}

On this subset, the best SAN operating points are Delta = 0.45 for
ResNet-50 (t* = 4, 44.5\% MAC saving, 7/8 clauses) and Delta = 0.55 for
ViT-small/d384 (t* = 4, \textasciitilde50\% MAC saving, 7/8 clauses).
The residual family benefits from a softer threshold because its
confidence grows more gradually across exits; the attention family needs
a sharper threshold to avoid firing on weak early heads. These values
are used as the starting points for the CIFAR-100 runs in Section4.4.

\hypertarget{larger-proxy-cifar-100}{%
\subsubsection{4.4 Larger proxy:
CIFAR-100}\label{larger-proxy-cifar-100}}

CIFAR-100 was staged on OrangeFS
(\texttt{/orangefs/training/sounio/datasets/cifar-100-python}) and two
Slurm jobs were submitted with thresholds adjusted to the harder,
100-class task. The first submission (jobs 8609/8610) failed because the
harness's CIFAR-100 loader used \texttt{encoding="latin1"}, under which
the pickle keys are strings rather than bytes; this was fixed to
\texttt{encoding="bytes"} and the runs were resubmitted as jobs 8611
(ResNet-50) and 8612 (ViT-small/d384).

On the same 4,000 train / 1,000 val split, CIFAR-100 is too hard for the
small epoch budget: neither family reaches its lowered tau, so
\texttt{t*\ =\ None} and the freeze-on-green rule never fires. The
result is reported as a negative result, not hidden.

\textbf{Table 4: CIFAR-100 small-subset pilot.}

\begin{longtable}[]{@{}
  >{\raggedright\arraybackslash}p{(\columnwidth - 14\tabcolsep) * \real{0.1250}}
  >{\raggedright\arraybackslash}p{(\columnwidth - 14\tabcolsep) * \real{0.1250}}
  >{\raggedright\arraybackslash}p{(\columnwidth - 14\tabcolsep) * \real{0.1250}}
  >{\raggedright\arraybackslash}p{(\columnwidth - 14\tabcolsep) * \real{0.1250}}
  >{\raggedright\arraybackslash}p{(\columnwidth - 14\tabcolsep) * \real{0.1250}}
  >{\raggedright\arraybackslash}p{(\columnwidth - 14\tabcolsep) * \real{0.1250}}
  >{\raggedright\arraybackslash}p{(\columnwidth - 14\tabcolsep) * \real{0.1250}}
  >{\raggedright\arraybackslash}p{(\columnwidth - 14\tabcolsep) * \real{0.1250}}@{}}
\toprule\noalign{}
\begin{minipage}[b]{\linewidth}\raggedright
family
\end{minipage} & \begin{minipage}[b]{\linewidth}\raggedright
tau
\end{minipage} & \begin{minipage}[b]{\linewidth}\raggedright
epochs
\end{minipage} & \begin{minipage}[b]{\linewidth}\raggedright
final acc
\end{minipage} & \begin{minipage}[b]{\linewidth}\raggedright
S\_m(SAN)
\end{minipage} & \begin{minipage}[b]{\linewidth}\raggedright
S\_m(Dense)
\end{minipage} & \begin{minipage}[b]{\linewidth}\raggedright
saving
\end{minipage} & \begin{minipage}[b]{\linewidth}\raggedright
verdict
\end{minipage} \\
\midrule\noalign{}
\endhead
\bottomrule\noalign{}
\endlastfoot
ResNet-50 & 0.20 & 8 & 0.126 & 266.9 TMAC & 270.0 TMAC & 1.2\% & L\_RED
(1/8) \\
ViT-small/d384 & 0.15 & 10 & 0.117 & 369.1 TMAC & 369.3 TMAC &
\textless0.1\% & L\_RED (infeasible) \\
\end{longtable}

The near-zero savings show that early exits are almost irrelevant when
the model is underfit: almost no sample becomes confident enough to
exit, so SAN runs essentially the full trunk. A usable CIFAR-100 result
would require a larger train split, more epochs, or a smaller model.
ImageNette-320/full remains a candidate real-image proxy if bandwidth
and storage permit.

\hypertarget{real-image-validation-on-imagenette2-160}{%
\subsubsection{4.5 Real-image validation on
ImageNette2-160}\label{real-image-validation-on-imagenette2-160}}

A SAN-ResNet-18 was trained on 4,000 ImageNette2-160 real photographs
(frozen ImageNet-1k backbone, layer4 + heads fine-tuned). Validation
confidences for 3,925 real images were exported to the U250 cohort
format and run on-target:

\begin{verbatim}
host_san_scan: dataset=val_imagenette n=3925 points=5 q_delta=18021 family=resnet
CARD_RESULT n=3925 catastrophes=241 flops_macs=4446713384960 wall=0.095ms (41.2Msamples/s kernel-only)
HOST_SAN_SCAN_PASS (val_imagenette)
\end{verbatim}

The card's histogram, catastrophe count, and FLOP total are bit-exact
against the Python golden model on real photographs.

\hypertarget{end-to-end-deployment-loop}{%
\subsubsection{4.6 End-to-end deployment
loop}\label{end-to-end-deployment-loop}}

Table 1 showed the scan kernel in isolation. We now report the first
measured pass through the full inference pipeline: a SAN-ResNet-18 trunk
running on the DL380 host CPU, exit confidences quantized to Q0.15,
packed into the same 512-bit beats used by \texttt{host\_san\_scan}, and
streamed to the U250 via the new \texttt{host\_san\_scan\_e2e} XRT host.
The orchestration script
\texttt{scripts/research/san\_fpga\_endtoend.py} measures every phase
and validates the card output bit-exactly against an independent Python
golden scan.

\textbf{Table 5: Host\textless-\textgreater card phase decomposition
(ImageNette2-160, n=3 925).}

\begin{longtable}[]{@{}
  >{\raggedright\arraybackslash}p{(\columnwidth - 4\tabcolsep) * \real{0.3333}}
  >{\raggedright\arraybackslash}p{(\columnwidth - 4\tabcolsep) * \real{0.3333}}
  >{\raggedright\arraybackslash}p{(\columnwidth - 4\tabcolsep) * \real{0.3333}}@{}}
\toprule\noalign{}
\begin{minipage}[b]{\linewidth}\raggedright
phase
\end{minipage} & \begin{minipage}[b]{\linewidth}\raggedright
time
\end{minipage} & \begin{minipage}[b]{\linewidth}\raggedright
note
\end{minipage} \\
\midrule\noalign{}
\endhead
\bottomrule\noalign{}
\endlastfoot
quantize + pack & \textasciitilde40 ms & Q0.15 floor + 512-bit beat
packing \\
xclbin setup & \textasciitilde135 ms & one-time \texttt{xclbin} load per
process (no PR) \\
DMA H2D & \textasciitilde0.12 ms & 62 848 bytes = 982 beats x 64
bytes \\
kernel & \textasciitilde0.66 ms & \textasciitilde6 Msamples/s
single-shot for this small cohort \\
DMA D2H & \textasciitilde0.15 ms & histogram + catastrophe count + MAC
total \\
\textbf{host\textless-\textgreater card total} &
\textbf{\textasciitilde136 ms} & first cohort; subsequent cohorts reuse
context \\
PyTorch forward (CPU) & \textasciitilde18.4 s & SAN-ResNet-18, 3 925
real images, DL380 CPU \\
\end{longtable}

The host\textless-\textgreater card total is the measured FPGA round
trip for one cohort; it is dominated by the one-time \texttt{xclbin}
load. After the load, additional cohorts can be enqueued without
reloading, so the per-cohort cost excluding forward is \textasciitilde1
ms. The PyTorch forward pass is CPU-only because the DL380 has no GPU
fitted; a GPU would reduce it proportionally. The packed cohort is 62
848 bytes = 982 beats x 64 bytes, so the DMA times correspond to
\textasciitilde470 MB/s effective PCIe bandwidth for the small transfer.
The script supports \texttt{-\/-mock-host} for CI or environments
without the U250; the mock run reproduces the same histogram,
catastrophe count, and FLOP total and confirms the Python packing
matches the C++ unpacking bit-exactly.

This closes the loop between training (Section4.2), confidence
extraction, and FPGA audit: the trunk runs on the host, the card decides
and meters, and the orchestration script verifies that the two agree
exactly.

\begin{center}\rule{0.5\linewidth}{0.5pt}\end{center}

\hypertarget{discussion}{%
\subsection{5 Discussion}\label{discussion}}

\hypertarget{what-the-numbers-mean}{%
\subsubsection{5.1 What the numbers mean}\label{what-the-numbers-mean}}

The U250 result shows that the SAN decision/metering path is deployable
as a small, fast, low-energy FPGA kernel. The kernel does not accelerate
the trunk; it accelerates the audit: for every incoming cohort, it
decides when each sample should have exited and counts the exact
computation that was executed. This is the operation a production SAN
deployment runs continuously.

The GPU result shows that the SAN training rule saves substantial
computation across architecture families on a real dataset, even at
CIFAR-10 scale. The savings would be larger at ImageNet scale because
the per-exit stage cost is much higher.

\hypertarget{limitations-stated-honestly}{%
\subsubsection{5.2 Limitations (stated
honestly)}\label{limitations-stated-honestly}}

\textbf{Small-subset training study.} The GPU study uses 4,000 train /
1,000 val images and low feasibility targets (tau = 0.34/0.251/0.165).
The reported accuracies (0.39, 0.26, 0.17) demonstrate the SAN training
machinery, not competitive CIFAR-10 performance. Claims are framed as
machinery validation and metered-MAC savings under identical accounting,
not as accuracy results.

\textbf{Partial meter convention.} The machine channel charges only
conv/linear MACs and attention token-mixing matmuls; biases, norms,
activations, softmax, residuals, pooling, host dispatch, and PCIe
transfers are unmetered. Reported percentages are relative savings under
this stated convention, not claims about total wall-time or total
energy.

\textbf{Full ImageNet-1k is unavailable in this environment.} It
requires -credentials and \textasciitilde150 GB of download; this node
has neither. ImageNette2-160 is the honest real-image proxy, and all
``ImageNet scale'' claims refer to (a) the real architecture FLOP
constants, (b) the 1.2M-sample stress cohort, or (c) explicit
extrapolation.

\textbf{U250 energy is board-level only.} We report the card's own power
sensor at 1 Hz, not a server-rack measurement. The incremental
\textasciitilde1.7 W is rounded to two significant figures;
three-decimal-nJ precision would require a higher-rate external meter.

\textbf{Patient channel is mixed.} ResNet-50 shows the disclosed
patient-channel tradeoff on this task: its integrated patient harm is
slightly above EarlyStop's because early stopping alone can freeze on a
luckier epoch. ViT-small/d384 and SAN-GPT-small satisfy the stricter L5
clause (SAN \textless= both baselines) in this run. The parent line
still reports an attention-family tradeoff at ImageNet scale, so the
small-proxy result should not be read as a general claim.

\textbf{Single bitstream, no DSE.} The work did not include
placement/routing exploration, multi-bitstream campaigns, or comparison
with HLS alternatives. The 135.2 MHz achieved clock is the honest result
of one Vitis build.

\textbf{No CPU baseline for the audit kernel.} Table 5 decomposes the
host\textless-\textgreater card phases, but we do not report a host-CPU
implementation of the same integer audit/metering kernel. The FPGA
speedup claim is therefore relative to the bus-limited theoretical peak
and to the application need (run the trunk on host/GPU and audit exits
on card), not relative to a measured CPU baseline.

\textbf{Table 2 omits the EarlyStop baseline.} The companion contract
evaluates SAN against both Dense and EarlyStop (L5), and the text
discusses the comparison, but the table itself reports only SAN and
Dense integrated machine burden. A future revision will add the
per-family EarlyStop column so the reader can see the savings against
both baselines in one view.

\hypertarget{future-work}{%
\subsubsection{5.3 Future work}\label{future-work}}

\begin{itemize}
\tightlist
\item
  Run the full pipeline on ImageNet-1k when credentials and storage
  become available.
\item
  Explore the threshold (Delta) and feasibility (tau) sensitivity
  surface across families.
\item
  Extend the kernel to multi-batch continuous streaming and measure host
  PCIe overhead in a server context.
\item
  Investigate sparsity and quantization ladders for the machine channel.
\item
  Add a measured host-CPU baseline for the integer audit/metering kernel
  to close the FPGA speedup claim.
\item
  Include the EarlyStop baseline column in the GPU training table for
  direct three-way comparison.
\end{itemize}

\begin{center}\rule{0.5\linewidth}{0.5pt}\end{center}

\hypertarget{conclusion}{%
\subsection{6 Conclusion}\label{conclusion}}

We have presented the first measured deployment of the SAN exit-audit /
FLOP-metering kernel on an FPGA, together with a GPU training study
across three architecture families. The U250 kernel runs at 511
Msamples/s on a 1.2M-sample stress cohort and approximately 3.3
nJ/sample board-level energy, with bit-exact correctness verified
against a golden model on synthetic, stress, and real-image cohorts. GPU
training saves 40.7--52.2\% of metered-MAC burden in a small-subset,
fast-convergence regime. We have stated the limitations plainly ---
small-subset accuracy, partial meter convention, no full ImageNet-1k,
board-level power only, ResNet-50 patient-channel tradeoff --- because
the honesty of the evidence is itself a contribution. The artifacts are
available for reproduction.

\begin{center}\rule{0.5\linewidth}{0.5pt}\end{center}

\hypertarget{ai-disclosure}{%
\subsection{7 AI Disclosure}\label{ai-disclosure}}

This draft was prepared under human direction, with mandatory
LLM-offload review per \texttt{.claude/AGENT\_OFFLOAD\_POLICY.md}. The
companion spec and contract were audited by math-review offload
(xai/Grok 4.3 and Z.AI GLM) and revised per their findings. No clinical
content. GAIDeT-ICMJE 2025.

\begin{center}\rule{0.5\linewidth}{0.5pt}\end{center}

\hypertarget{references}{%
\subsection{References}\label{references}}

{[}1{]} Teerapittayanon, S., McDanel, B., \& Kung, H. T. (2016).
BranchyNet: Fast inference via early exiting from deep neural networks.
\emph{ICPR}.

{[}2{]} Kaya, Y., Hong, S., \& Dumitras, T. (2019). Shallow-Deep
Networks: Understanding and mitigating network overthinking.
\emph{ICML}.

{[}3{]} Huang, G. (2018). Multi-scale dense networks for resource
efficient image classification. \emph{ICLR}.

{[}4{]} Sounio repository, suffering-aware neural network line:
\texttt{docs/research/suffering\_aware\_architecture\_spec\_2026-07-28.md}
and successors.

{[}5{]} Chamon, L. F., \& Ribeiro, A. (2020). Probably approximately
correct constrained learning. \emph{NeurIPS}.

{[}6{]} Chamon, L. F., Paternain, S., Calvo-Fullana, M., \& Ribeiro, A.
(2023). Constrained learning with non-convex losses. \emph{IEEE
Transactions on Information Theory}.

{[}7{]} Cotter, A., Jiang, H., Gupta, M., Wang, S., Narayan, T., You,
S., \& Sridharan, D. (2019). Optimization with non-differentiable
constraints with applications to fairness, recall, churn, and other
goals. \emph{JMLR}.

{[}8{]} Manheim, D., \& Garrabrant, S. (2018). Categorizing variants of
Goodhart's law. arXiv:1803.04585.

\begin{center}\rule{0.5\linewidth}{0.5pt}\end{center}

\hypertarget{a-reproduction}{%
\subsection{A Reproduction}\label{a-reproduction}}

\begin{Shaded}
\begin{Highlighting}[]
\CommentTok{\# CPU contract + spec/outline consistency (\textasciitilde{}3 min on CPU)}
\FunctionTok{bash}\NormalTok{ scripts/ci/san\_imagenet\_fpga\_dl380\_gate.sh}
\CommentTok{\# expect: SAN\_IMAGENET\_FPGA\_DL380\_VERDICT I\_GREEN (8/8 clauses PASS)}

\CommentTok{\# GPU training (requires CUDA)}
\VariableTok{SAN\_LARGE\_DEVICE}\OperatorTok{=}\NormalTok{cuda }\VariableTok{SAN\_LARGE\_ONLY}\OperatorTok{=}\NormalTok{resnet50 }\DataTypeTok{\textbackslash{}}
  \ExtensionTok{.venv/bin/python}\NormalTok{ scripts/research/suffering\_aware\_large\_architecture.py}

\CommentTok{\# U250 host scan (requires Xilinx XRT + bitstream on DL380)}
\ExtensionTok{./host\_san\_scan} \AttributeTok{{-}x}\NormalTok{ krnl\_san\_scan.hw.xclbin }\AttributeTok{{-}d}\NormalTok{ val\_resnet}
\end{Highlighting}
\end{Shaded}

\begin{center}\rule{0.5\linewidth}{0.5pt}\end{center}

\hypertarget{b-companion-contract-clauses-l1l8}{%
\subsection{B Companion contract clauses
(L1--L8)}\label{b-companion-contract-clauses-l1l8}}

The GPU harness reports per-family verdicts against clauses L1--L8 of
\texttt{scripts/research/suffering\_aware\_large\_architecture.py}:

\begin{longtable}[]{@{}
  >{\raggedright\arraybackslash}p{(\columnwidth - 4\tabcolsep) * \real{0.3333}}
  >{\raggedright\arraybackslash}p{(\columnwidth - 4\tabcolsep) * \real{0.3333}}
  >{\raggedright\arraybackslash}p{(\columnwidth - 4\tabcolsep) * \real{0.3333}}@{}}
\toprule\noalign{}
\begin{minipage}[b]{\linewidth}\raggedright
clause
\end{minipage} & \begin{minipage}[b]{\linewidth}\raggedright
predicate
\end{minipage} & \begin{minipage}[b]{\linewidth}\raggedright
threshold / note
\end{minipage} \\
\midrule\noalign{}
\endhead
\bottomrule\noalign{}
\endlastfoot
L1 & Metering conservation: gated-off stages charge 0; metered MACs ==
independent manual accounting & equality \\
L2 & Feasibility: SAN reaches val acc \textgreater= tau within budget &
tau per family \\
L3 & Anti-Goodhart: 101-weight compassion grid never selects an
infeasible checkpoint & abstainer \& probe infeasible \\
L4 & Necessary/gratuitous separation: SAN gratuitous machine burden = 0
& equality \\
L5 & Suffering bounds: SAN machine burden \textless{} Dense and
\textless= EarlyStop; patient harm \textless= both baselines & within
small tolerance \\
L6 & Exits are real: val exit fraction at t* \textgreater{} 0.07 &
0.07 \\
L7 & Patient channel first-class: harm matrix off-diagonal max
\textgreater= 3x min; SAN peak harm \textless= baselines & 3x
asymmetry \\
L8 & Anti-shortcut: a spurious-feature probe that overfits train but
fails val is rejected by the gate & train \textgreater{} tau, val
\textless{} tau, never selected \\
\end{longtable}

The FPGA / ImageNet-scale contract uses an analogous I1--I8 clause set
defined in \texttt{scripts/research/san\_imagenet\_fpga\_dl380.py} and
enforced by \texttt{scripts/ci/san\_imagenet\_fpga\_dl380\_gate.sh}.

\end{document}
